\uuid{Ob2N}
\titre{Étude de la différentiabilité de $f(x,y)=\dfrac{x^2 y}{x^2+y^2}$}
\niveau{L2}
\module{Analyse}
\chapitre{Fonctions de plusieurs variables}
\sousChapitre{Continuité et différentiabilité}
\theme{Classe $\mathcal{C}^1$, différentiabilité, coordonnées polaires, limite dépendant de $\theta$}
\auteur{Maxime Nguyen}
\datecreate{2026-03-24}
\organisation{amscc}
\difficulte{4}

\contenu{
\texte{
  On considère la fonction $f$ définie sur $\mathbb{R}^2$ par
  \[
    f(x,y) = \frac{x^2 y}{x^2+y^2} \;\text{ si } (x,y) \neq (0,0), \qquad f(0,0) = 0.
  \]
}
\begin{enumerate}
  \item \question{La fonction $f$ est-elle $\mathcal{C}^1$ en $(0,0)$ ?}
    \reponse{
      Pour $(x,y) \neq (0,0)$ :
      \[
        \frac{\partial f}{\partial x}(x,y) = \frac{2xy^3}{(x^2+y^2)^2}, \qquad \frac{\partial f}{\partial y}(x,y) = \frac{x^2(x^2-y^2)}{(x^2+y^2)^2}.
      \]
      En coordonnées polaires :
      \[
        \frac{\partial f}{\partial x}(x,y) = 2\cos\theta\sin^3\theta,
      \]
      qui dépend de $\theta$ et n'a donc pas de limite en $(0,0)$. Ainsi $\dfrac{\partial f}{\partial x}$ n'est pas continue en $(0,0)$, et $f$ n'est pas $\mathcal{C}^1$ en $(0,0)$.
    }

  \item \question{La fonction $f$ est-elle différentiable en $(0,0)$ ?}
    \reponse{
      On calcule $\dfrac{\partial f}{\partial x}(0,0) = \dfrac{\partial f}{\partial y}(0,0) = 0$. Pour la différentiabilité, on étudie :
      \[
        \frac{f(x,y)}{\sqrt{x^2+y^2}} = \frac{r\cos^2\theta\sin\theta}{r} = \cos^2\theta\sin\theta,
      \]
      qui dépend de $\theta$. Donc $f$ n'est pas différentiable en $(0,0)$.
    }
\end{enumerate}
}
